\begin{lemma}[The envelopes are genuine suprema and infima]
Let $F$ be a finite-dimensional real inner-product space, let $I\subseteq\R$,
let $f$ be a real function on the space-time $\ell^2(F\times\R)$ which is
continuous on $\Omega_I=\{(y,\tau):\tau\in I\}$ and has subquadratic spatial
growth on compact time subintervals of $I$ in the sense of
\noderef{SubquadraticGrowth}. Let $z=(x,t)$ with $t\in I$. Then:
\begin{enumerate}
\item for every $\ell>0$ with $t+\ell\in I$, the forward envelope
$\Penv_\ell f(z)$ of \noderef{QuadEnvP} dominates every competitor and is the
least such bound:
\[
  f(y,t+\ell)-\frac{|y-x|^2}{2\ell}\le\Penv_\ell f(z)\quad(y\in F),
  \qquad
  \Bigl(\forall y\ \ f(y,t+\ell)-\frac{|y-x|^2}{2\ell}\le b\Bigr)
     \Longrightarrow \Penv_\ell f(z)\le b;
\]
\item for every $\ell>0$ with $t-\ell\in I$, the backward envelope
$\Menv_\ell f(z)$ of \noderef{QuadEnvM} is dominated by every competitor and is
the greatest such bound:
\[
  \Menv_\ell f(z)\le f(y,t-\ell)+\frac{|y-x|^2}{2\ell}\quad(y\in F),
  \qquad
  \Bigl(\forall y\ \ b\le f(y,t-\ell)+\frac{|y-x|^2}{2\ell}\Bigr)
     \Longrightarrow b\le\Menv_\ell f(z);
\]
\item for every $s>0$, the barycentric envelope $\Cenv_sf(z)$ of
\noderef{BaryEnv} dominates the value of every finite horizontal probability
configuration $(\theta_i,y_i)$ with mean $x$ in the sense of
\noderef{IsBaryConfig}, with radius $\rho$ the root mean square radius of
\noderef{BaryRadius}, and is the least such bound:
\[
  \sum_i\theta_if(y_i,t)-\frac{\rho^2}{2s}\le\Cenv_sf(z),
  \qquad
  \Bigl(\forall\ \text{configurations}\ \
     \sum_i\theta_if(y_i,t)-\frac{\rho^2}{2s}\le b\Bigr)
     \Longrightarrow \Cenv_sf(z)\le b .
\]
\end{enumerate}
\end{lemma}
\begin{proof}
The three envelopes of \noderef{QuadEnvP}, \noderef{QuadEnvM} and
\noderef{BaryEnv} are, for the parameters considered here, the supremum, the
infimum and the supremum of the three sets
\begin{align*}
  \mathcal P&=\Bigl\{f(y,t+\ell)-\frac{|y-x|^2}{2\ell}:y\in F\Bigr\},\\
  \mathcal M&=\Bigl\{f(y,t-\ell)+\frac{|y-x|^2}{2\ell}:y\in F\Bigr\},\\
  \mathcal C&=\Bigl\{\sum_i\theta_if(y_i,t)-\frac{\rho^2}{2s}:
      (\theta_i,y_i)\ \text{a configuration with mean }x
      \text{ and radius }\rho\Bigr\},
\end{align*}
respectively. By the preamble's convention a supremum is the least upper bound
and an infimum the greatest lower bound exactly when the set is nonempty and
bounded on the relevant side, and otherwise the value is $0$ and carries no
information. So each of the three assertions is precisely the statement that the
corresponding set is nonempty and bounded above, respectively below; once that
is known, the two displayed properties are the defining properties of the least
upper bound and the greatest lower bound. We prove the three boundedness
statements. Note that in each case the time argument of $f$ lies in $I$, by the
hypotheses $t+\ell\in I$, $t-\ell\in I$ and $t\in I$, so $f$ is continuous at
the points used.

\emph{A remark on why this is not a restatement of attainment.} The assertions
above quantify over \emph{all} competitors. A statement that the supremum is
\emph{attained} at some single $y$ --- as in the node
\texttt{QuadEnvAttained}, part (1), or in the uniqueness clause of the node
\texttt{DPSpatialRegularity} --- names one
optimizer and does not, under the preamble's convention, entail that the value
$\Penv_\ell f(z)$ dominates the competitor at any other point: if $\mathcal P$
were unbounded above then $\Penv_\ell f(z)$ would be $0$, which may coincide
with the value at some particular $y$ while failing the first display. Nor does
boundedness entail attainment, since a bounded set need not contain its
supremum. The two statements are therefore independent, and only the present
one licenses a comparison of the envelope with a chosen competitor.

\emph{Part (1).} The set $\mathcal P$ is nonempty: the choice $y=x$ contributes
$f(x,t+\ell)$. For the upper bound, apply \noderef{SubquadraticGrowth} with the
compact set $J=\{t+\ell\}\subseteq I$ and the accuracy
$\varepsilon=\dfrac{1}{4\ell}>0$: there is $M>0$ with
\begin{equation}\label{eq:ecb-far}
  |f(y,t+\ell)|\le\frac{|y|^2}{4\ell}\qquad(|y|\ge M).
\end{equation}
The function $f(\cdot,t+\ell)$ is continuous on the closed ball
$\overline B_M(0)$ of $F$, because every point $(y,t+\ell)$ with $y$ in that
ball lies in $\Omega_I$; that ball is compact, so there is $C_M<\infty$ with
\begin{equation}\label{eq:ecb-near}
  |f(y,t+\ell)|\le C_M\qquad(|y|\le M).
\end{equation}
Enlarging $C_M$ we may assume $C_M\ge0$. From
$|y|^2\le2|y-x|^2+2|x|^2$ we get
\begin{equation}\label{eq:ecb-quad}
  |y-x|^2\ge\frac{|y|^2}{2}-|x|^2\qquad(y\in F).
\end{equation}
Now let $y\in F$. If $|y|\ge M$ then by \eqref{eq:ecb-far} and
\eqref{eq:ecb-quad}
\[
  f(y,t+\ell)-\frac{|y-x|^2}{2\ell}
   \le\frac{|y|^2}{4\ell}
       -\frac{1}{2\ell}\Bigl(\frac{|y|^2}{2}-|x|^2\Bigr)
   =\frac{|x|^2}{2\ell},
\]
since $\dfrac{|y|^2}{4\ell}-\dfrac{|y|^2}{4\ell}=0$. If $|y|\le M$ then by
\eqref{eq:ecb-near} and $|y-x|^2\ge0$,
\[
  f(y,t+\ell)-\frac{|y-x|^2}{2\ell}\le C_M .
\]
In both cases the value is at most $C_M+\dfrac{|x|^2}{2\ell}$, a bound
independent of $y$. Hence $\mathcal P$ is bounded above, which proves part (1).

\emph{Part (2).} Apply part (1) to the function $-f$ and the time $t-\ell$. The
function $-f$ is continuous on $\Omega_I$ and satisfies the hypothesis of
\noderef{SubquadraticGrowth} with the same data, because $|-f|=|f|$. Part (1)'s
boundedness argument, run with $t-\ell\in I$ in place of $t+\ell$, shows that
\[
  \Bigl\{-f(y,t-\ell)-\frac{|y-x|^2}{2\ell}:y\in F\Bigr\}
\]
is nonempty and bounded above, by $C'_M+\dfrac{|x|^2}{2\ell}$ where $C'_M$
bounds $|f(\cdot,t-\ell)|$ on $\overline B_M(0)$ for the $M$ supplied by
\noderef{SubquadraticGrowth} at $J=\{t-\ell\}$ and
$\varepsilon=1/(4\ell)$. Multiplying by $-1$ reverses the inequalities and shows
that $\mathcal M$ is nonempty and bounded below, by
$-C'_M-\dfrac{|x|^2}{2\ell}$. This proves part (2).

\emph{Part (3).} The set $\mathcal C$ is nonempty: by \noderef{IsBaryConfig} the
one-point configuration $n=1$, $\theta_1=1$, $y_1=x$ is admissible, since its
weight is positive, the weights sum to $1$ and its weighted mean is $x$; by
\noderef{BaryRadius} its radius is $\rho=0$, so it contributes $f(x,t)$. For the
upper bound, apply \noderef{SubquadraticGrowth} with the compact set
$J=\{t\}\subseteq I$ and the accuracy $\varepsilon=\dfrac{1}{4s}>0$ to get
$M>0$ with
\begin{equation}\label{eq:ecb-bary-far}
  |f(y,t)|\le\frac{|y|^2}{4s}\qquad(|y|\ge M),
\end{equation}
and let $C_M\ge0$ bound the continuous $|f(\cdot,t)|$ on the compact ball
$\overline B_M(0)$. Let $(\theta_i,y_i)_{i\le n}$ be a configuration with mean
$x$ and radius $\rho$, so that $\theta_i>0$, $\sum_i\theta_i=1$ and
$\sum_i\theta_iy_i=x$ by \noderef{IsBaryConfig}, and
$\rho^2=\sum_i\theta_i|y_i-x|^2$ by \noderef{BaryRadius}. Expanding the squared
norms and using these two identities,
\[
  \sum_i\theta_i|y_i-x|^2
   =\sum_i\theta_i|y_i|^2
      -2\Bigl\langle\sum_i\theta_iy_i,\,x\Bigr\rangle
      +\Bigl(\sum_i\theta_i\Bigr)|x|^2
   =\sum_i\theta_i|y_i|^2-|x|^2,
\]
that is
\begin{equation}\label{eq:ecb-variance}
  \sum_i\theta_i|y_i|^2=\rho^2+|x|^2 .
\end{equation}
For each index $i$, \eqref{eq:ecb-bary-far} in the case $|y_i|\ge M$ and the
bound $C_M$ in the case $|y_i|\le M$ give
\[
  f(y_i,t)\le C_M+\frac{|y_i|^2}{4s},
\]
since both terms on the right are nonnegative. Multiplying by $\theta_i>0$,
summing, and using $\sum_i\theta_i=1$ together with \eqref{eq:ecb-variance},
\[
  \sum_i\theta_if(y_i,t)\le C_M+\frac{\rho^2+|x|^2}{4s}.
\]
Therefore the configuration value satisfies
\[
  \sum_i\theta_if(y_i,t)-\frac{\rho^2}{2s}
   \le C_M+\frac{|x|^2}{4s}+\rho^2\Bigl(\frac{1}{4s}-\frac{1}{2s}\Bigr)
   \le C_M+\frac{|x|^2}{4s},
\]
because $\dfrac1{4s}-\dfrac1{2s}=-\dfrac1{4s}<0$ and $\rho^2\ge0$. This bound is
independent of the configuration, so $\mathcal C$ is bounded above, which proves
part (3).
\end{proof}
